% Bagian: incompleteness
% Bab: theories-computability
% Subbagian: q-is-ce

\documentclass[../../../include/open-logic-section]{subfiles}

\begin{document}

\olfileid[id]{inc}{tcp}{qce}

\olsection{$\Th{Q}$ \printtoken{S}{c.e.}-Lengkap}


\begin{thm}
$\Th{Q}$ bersifat !!{c.e.}, tetapi tidak dapat diputuskan. Bahkan,
teori itu merupakan himpunan !!{c.e.} lengkap.
\end{thm}

\begin{proof}
Mudah dilihat bahwa $\Th{Q}$ bersifat !!{c.e.}, sebab teori itu merupakan
himpunan (kode) kalimat~$y$ sedemikian sehingga terdapat bukti~$x$ atas~$y$
dalam~$\Th{Q}$:
\[
Q = \Setabs{y}{\lexists[x][\Prf[\Th{Q}](x,y)]}.
\]
Namun, kita mengetahui bahwa $\Prf[\Th{Q}](x,y)$ dapat dihitung (bahkan
rekursif primitif), dan setiap himpunan yang dapat ditulis dalam bentuk di
atas bersifat c.e.

Pernyataan bahwa teori itu merupakan himpunan c.e. lengkap ekuivalen dengan
pernyataan bahwa $K \leq_m Q$, dengan $K = \Setabs{x}{!A_x(x) \downarrow}$.
Jadi, mari kita tunjukkan bahwa $K$ tereduksi ke~$\Th{Q}$. Karena predikat
Kleene~$T(e,x,s)$ bersifat rekursif primitif, predikat itu dapat
direpresentasikan dalam~$\Th{Q}$, misalnya oleh~$!A_T$. Maka untuk setiap~$x$,
kita mempunyai
\begin{align*}
x \in K & \lif \lexists[s][T(x,x,s)] \\
& \lif \lexists[s][(\Th{Q} \Proves !A_T(\num x,\num x, \num s))]
  \\
& \lif \Th{Q} \Proves \lexists[s][!A_T(\num x, \num x, s)].
\end{align*}
Sebaliknya, jika $\Th{Q} \Proves \lexists[s][!A_T(\num x, \num x, s)]$,
maka sesungguhnya, untuk suatu bilangan asli~$n$, formula $!A_T(\num x,
\num x, \num n)$ harus benar. Sekarang, jika $T(x,x,n)$ salah, $\Th{Q}$
akan membuktikan $\lnot !A_T(\num x, \num x, \num n)$, sebab $!A_T$
merepresentasikan~$T$. Namun, dengan demikian $\Th{Q}$ membuktikan suatu
formula yang salah, dan ini merupakan kontradiksi. Jadi $T(x,x,n)$ harus
benar, yang mengimplikasikan $!A_x(x) \downarrow$.

Ringkasnya, untuk setiap~$x$, $x$ berada dalam~$K$ jika dan hanya jika
$\Th{Q}$ membuktikan $\lexists[s][T(\num x,\num x,s)]$. Jadi fungsi~$f$
yang memetakan~$x$ ke (kode) kalimat $\lexists[s][T(\num x,
  \num x, s)]$ merupakan reduksi dari~$K$ ke~$\Th{Q}$.
\end{proof}

\end{document}
